\section{Roofline analysis}
\label{sec:roofline}

Figure \ref{fig:roofline} places the kernels against the ceilings. The bandwidth
ceiling is measured on this GPU. The compute roofs are the hardware's, computed as
48 SMs times the SM clock times 512 FLOP per cycle for FP16 inputs with FP32
accumulate, and 256 FLOP per cycle for FP32.

That distinction is the whole point of the chapter, and it is a correction. An
earlier version of this profiler set the roof to a cuBLAS measurement, which turns
every "percent of roof" into "percent of cuBLAS" drawn on log axes and makes the
chart look like evidence while carrying none. cuBLAS is drawn instead as a dashed
attainable line, and it reaches 97.2 percent of the tensor roof, which is itself
the best available evidence that 512 FLOP per cycle is the right rate for this
part.

The GEMM variants sit far to the right of both ridges, in the compute bound region,
with the top kernel at 84.1 percent of the tensor roof. GEMV sits at 0.5 FLOP per
byte on the bandwidth diagonal in the memory bound region. Those two points at
opposite ends of the intensity axis are the evidence behind the compute bound claim
and behind the tile choice.

\begin{figure}[htbp]
\centering
\includegraphics[width=0.95\textwidth]{roofline.pdf}
\caption{Roofline on the RTX 5070. The bandwidth ceiling is measured; the compute
roofs are the hardware's, and cuBLAS is a dashed attainable line rather than a
ceiling. \rooflinecaption}
\label{fig:roofline}
\end{figure}

The chart is built to take four families rather than one. The profiler emits an
operating point per variant and the plotter draws whichever points the CSV carries,
so the SpMV, FFT and scan points join the chart when their sweep rows land without
a code change. The caption above names which families are on the chart today and
which are pending, and it is generated rather than written, so it cannot claim a
family that is not there.

The roofline also settles one epilogue fusion decision, which is the smallest and
cleanest use of the model in the project. A standalone bias plus ReLU pass has an
arithmetic intensity of 0.167 FLOP per byte and is therefore firmly memory bound.
Folding it into a compute bound GEMM removes a pure bandwidth pass rather than
saving arithmetic, and the measured saving on the fused shape is 8.4 percent. That
is the fusion argument in miniature: the win is in the pass not taken, not in the
FLOP not executed~\cite{dao2022}.
